% iceCoder — arXiv-style short paper (compile with pdflatex or xelatex)
% Suggested categories: cs.SE, cs.AI
\documentclass[11pt,a4paper]{article}
\usepackage[margin=1in]{geometry}
\usepackage{amsmath,amssymb}
\usepackage{graphicx}
\usepackage{hyperref}
\usepackage{booktabs}
\usepackage{enumitem}

\title{iceCoder: A Selective Dual-Mode Runtime Supervisor\\
for Long-Horizon Tool-Using Coding Agents}

\author{
  \textbf{Anonymous Authors}\\
  \textit{iceCoder Open Source Project}\\
  \texttt{https://github.com/iceCoder-project}
}

\date{}

\begin{document}
\maketitle

\begin{abstract}
Autonomous coding agents that invoke tools over dozens of rounds frequently exhibit \emph{goal drift}, \emph{tool loops}, and \emph{stalling} without producing verifiable deliverables.
Productized assistants optimize model calls and chat UX, but offer limited, opaque runtime governance when sessions compress, crash, or exhaust context.
We present \textbf{iceCoder}, a self-hosted \emph{runtime governance layer} that wraps a standard tool-using LLM harness with three mechanisms: (i) explicit \textbf{deviation signals} observed each round; (ii) a \textbf{selective dual-mode} policy that defaults to free execution and escalates to forced constraints only on critical task domains; and (iii) \textbf{checkpoint-backed recovery} that persists task state beyond chat transcripts.
The design is formalized as a three-tier stack (policy tier L0, execution mode L1, runtime supervisor L2) with architectural axioms that centralize mode transitions and correction injection.
On two multi-file repair benchmarks under controlled conditions (same model \texttt{minimax-m2.5}, blind LLM-as-judge scoring), iceCoder matches or exceeds a Claude Code baseline on objective success while improving composite quality scores (+3 on a 100-point rubric in both tasks).
The system ships as TypeScript/Node.js with 1{,}165 unit tests and deployable via CLI, Web, WebSocket, and MCP.
\end{abstract}

\section{Introduction}
\label{sec:intro}

Large language models (LLMs) embedded in \emph{agent harnesses}---loops that alternate between generation and tool execution---can repair repositories, run tests, and iterate for twenty or more rounds.
In practice, failures are often \emph{runtime} failures: the model explores read-only tools during an edit task, retries identical failing commands, or edits the same file repeatedly without passing verification.
Users discover these problems only after wasted tokens and time.

iceCoder targets this gap.
It is \textbf{not} an IDE replacement; it is an attachable layer that can run beneath or beside existing agent stacks (CLI, Web UI, WebSocket API, optional MCP tools).
Our thesis is simple: \emph{trust the model by default, supervise selectively, recover from state on disk}.

Contributions:
\begin{enumerate}[nosep]
  \item A \textbf{three-tier supervision architecture} (L0 global policy, L1 free/forced execution mode, L2 passive observation and recovery) with enforceable axioms (single correction port, single mode authority, bounded correction budget).
  \item A \textbf{TaskGraph} integration that is the sole structured execution context injected into prompts for critical intents, plus domain gating that keeps Q\&A and inspection in free mode.
  \item \textbf{CheckpointEngine v2}, merging legacy task checkpoints with additive \texttt{runtimeV2} telemetry for post-compaction and process restart.
  \item An initial \textbf{controlled benchmark protocol} (same-model, blind judge, gate + six-dimensional rubric) with two published multi-file engineering tasks.
\end{enumerate}

\section{Background and Related Work}
\label{sec:related}

\textbf{Tool-using agents.}
Systems such as ReAct, code interpreters, and IDE-integrated agents (e.g., Claude Code, Cursor Composer, Codex-class CLIs) share a common pattern: LLM $\rightarrow$ tool calls $\rightarrow$ observation $\rightarrow$ repeat.
iceCoder assumes this pattern and adds governance around it.

\textbf{Planning vs.\ supervision.}
Classical workflow engines and explicit planners trade flexibility for control.
iceCoder explicitly \emph{rejects} replacing the LLM with a rigid workflow; the TaskGraph supplies \emph{context and validation}, not a hard orchestrator.
The dual-mode supervisor is a \textbf{runtime kernel}, not a planner.

\textbf{Recovery and memory.}
Session compaction and checkpointing appear in production agents, but often as opaque product features.
iceCoder makes recovery a first-class file artifact (\texttt{\{sessionId\}.checkpoint.json}) and couples it with runtime recovery prompts after compaction.

\textbf{Evaluation.}
Prior agent benchmarks emphasize single-shot coding or human preference.
We adopt a \textbf{dual-track} score: an objective gate (tests, build, file-scope compliance) plus blind LLM-as-judge dimensions (correctness, minimal change, verification discipline).

\section{System Overview}
\label{sec:system}

\subsection{Harness loop}

At the core is a \textbf{Harness} state machine: initialize task and repository ledgers; each round optionally compact context, inject runtime/repo state, recall file-based memories, call the LLM with tools, apply permissions, execute tools, update state, and enforce a verification gate before stop.
Protections include no-tool recovery for executable tasks, repeat-failure signatures, consecutive-failure circuit breaking, and an isolated \textbf{sub-agent} for read-only exploration (returning summaries instead of polluting the main context).

\subsection{Three-tier dual-mode supervision}
\label{sec:dualmode}

\begin{figure}[ht]
\centering
\fbox{\parbox{0.92\linewidth}{
\small
\textbf{L0 Policy} (\texttt{off} / \texttt{adaptive} / \texttt{strict}) --- user-facing tier\\
$\downarrow$\\
\textbf{L1 Execution Mode} (\texttt{free} $\leftrightarrow$ \texttt{forced}) --- ToolGate, TaskGraph, branch budget\\
$\downarrow$\\
\textbf{L2 Runtime Supervisor} --- PassiveObserver, GoalDriftDetector, RecoverySupervisor, CorrectionPort
}}
\caption{Supervision stack (V1.3.7). L0 configures whether automatic escalation is enabled; only L1's ModeDecisionEngine may switch execution mode.}
\label{fig:stack}
\end{figure}

\textbf{L0} selects global policy: disable supervision, adapt by risk (default), or enforce strict constraints (including first-round graph initialization).

\textbf{L1} implements \texttt{free} versus \texttt{forced} execution.
In \texttt{free} mode the agent explores with lighter gates; in \texttt{forced} mode stricter tool gates, structured step context from TaskGraph, and verification expectations apply.
Escalation is driven by \emph{signals}, not by parsing user keywords as triggers.

\textbf{L2} observes without directly mutating messages.
\texttt{PassiveObserver} accumulates deviation signals each tool round:
\begin{itemize}[nosep]
  \item \texttt{no\_progress} --- consecutive read-only or all-failed rounds;
  \item \texttt{tool\_repeat\_fail} --- repeated identical failing tool signatures;
  \item \texttt{file\_loop} --- excessive edits to one file (strict tier);
  \item \texttt{goal\_drift} --- heuristic misalignment between intent, tools, files touched, and assistant text (GoalDriftDetector).
\end{itemize}
\texttt{RecoverySupervisor} may \emph{take over} with bounded corrections via a single \texttt{CorrectionPort}, respecting recovery budgets and minimum forced dwell time (axiom I10) to prevent mode flicker.

Key axioms (abbreviated): corrections only through CorrectionPort (I1); ToolGate blocks are mandatory (I2); adaptive free segments skip graph init by default (I3); correction hints are budgeted (I4); only ModeDecisionEngine switches execution mode (I5); environment flags do not leak into per-task logic (I6).

\subsection{TaskGraph as sole structured context}

For critical intents (\texttt{edit}, \texttt{debug}, \texttt{refactor}, \texttt{test}, etc.), \texttt{GraphExecutor} builds a multi-step TaskGraph and injects \emph{only} the current node context into prompts.
Non-critical intents (\texttt{question}, \texttt{inspect}, \texttt{explain}) bypass graph initialization via \textbf{TaskDomainGate}, preserving free-mode behavior for low-risk interactions.

\subsection{CheckpointEngine v2}

Checkpoints store authoritative task fields for resume plus an additive \texttt{runtimeV2} object: recent tools, failures, recovery signals, branch-budget snapshots.
Writes are serialized to avoid torn JSON.
Triggers align with long-session milestones (tool steps, verification, compaction).

\subsection{File-based memory}

A separate file-memory subsystem provides long-term recall, session notes for compaction survival, eviction/dream consolidation, and security scanning---integrated at pre-LLM recall and post-tool extraction hooks.

\section{Evaluation}
\label{sec:eval}

\subsection{Protocol}

We use a \textbf{controlled comparison} protocol (JUDGE\_RUBRIC v0.1):
\begin{itemize}[nosep]
  \item Fixed model: \texttt{minimax-m2.5} across platforms;
  \item Identical prompts and sandbox repositories;
  \item Objective \textbf{Gate} (0--40): tests, build, file-scope compliance, no secret leakage;
  \item Blind \textbf{Judge} (0--60): six dimensions including requirement coverage, correctness, code quality, minimal change, verification discipline, and delivery notes;
  \item Composite $=$ Gate $+$ Judge; tiers S/A/B/C from composite thresholds.
\end{itemize}

iceCoder runs with \texttt{adaptive} supervisor mode unless noted.

\subsection{Tasks and results}

\textbf{Task 1: multi-file order pipeline.}
A TypeScript order-processing service with inventory, pricing, payment retry, and pipeline rollback bugs.
Both iceCoder and Claude Code achieved 9/9 tests and full gate points; iceCoder scored composite \textbf{86} vs.\ \textbf{83}, with advantages on correctness (transient retry semantics) and minimal change (no out-of-scope session artifacts).

\textbf{Task 2: saga warehouse reconciliation.}
A saga/event-sourcing exercise requiring reconciliation engine, compensation, and timeout monitoring.
Both platforms reached 15/15 tests; iceCoder scored \textbf{88} vs.\ \textbf{85}, with full gate credit (no auxiliary directory pollution) and comparable implementation depth.

Table~\ref{tab:results} summarizes published runs (May 2026).

\begin{table}[ht]
\centering
\caption{Benchmark summary (same model, blind judge). SR = objective success rate.}
\label{tab:results}
\begin{tabular}{lccc}
\toprule
Task & Platform & SR & Composite \\
\midrule
Order pipeline & iceCoder & \checkmark & 86 (A) \\
Order pipeline & Claude Code & \checkmark & 83 (A) \\
Saga reconciliation & iceCoder & \checkmark & 88 (A) \\
Saga reconciliation & Claude Code & \checkmark & 85 (A) \\
\bottomrule
\end{tabular}
\end{table}

\subsection{Regression assurance}

The open-source tree includes \textbf{93} Vitest files and \textbf{1{,}165} cases covering harness resilience, supervisor bridges, TaskGraph persistence, memory lifecycle, and web API routes---treated as the engineering source of truth alongside TypeScript type-checking.

\section{Discussion and Limitations}
\label{sec:limits}

\textbf{Scope.}
Results reflect two multi-file repair tasks and one model family; broader task suites (14-task rubric design) and cross-model studies remain future work.

\textbf{Goal drift detection.}
V1 uses lightweight heuristics (tool/intent consistency, file-touch alignment, bigram Jaccard on assistant text, progress factors); LLM-based gray-zone arbitration is specified but not the default path.

\textbf{Eval harness.}
\texttt{npm run eval:agent} remains a metric skeleton without a fully automated scoring runner in-tree.

\textbf{Deployment.}
iceCoder is self-hosted; operators must supply LLM provider configuration.
Supervision adds latency only on escalation paths; shadow mode (\texttt{ICE\_SUPERVISOR\_SHADOW=1}) supports observe-only rollout.

\section{Conclusion}
\label{sec:conclusion}

Long-horizon coding agents need explicit runtime governance: detect drift early, constrain execution selectively, and recover state beyond chat logs.
iceCoder instantiates this as a three-tier supervisor with architectural axioms, TaskGraph-gated structure, and checkpoint v2 resilience.
Initial controlled benchmarks suggest that governance can improve quality scores without sacrificing objective success under the same frontier model.
We release the implementation and benchmark reports to support reproducible agent-runtime research.

\section*{Acknowledgments}
We thank contributors to the iceCoder open-source repository and benchmark sandbox maintainers.

\begin{thebibliography}{9}

\bibitem{react}
S.~Yao et al.,
``ReAct: Synergizing Reasoning and Acting in Language Models,''
ICLR, 2023.

\bibitem{swelab}
C.~Jimenez et al.,
``SWE-bench: Can Language Models Resolve Real-World GitHub Issues?''
ICLR, 2024.

\bibitem{toolformer}
T.~Schick et al.,
``Toolformer: Language Models Can Teach Themselves to Use Tools,''
NeurIPS, 2023.

\bibitem{claude-code}
Anthropic,
``Claude Code --- agentic coding tool,''
product documentation, 2025--2026.

\bibitem{cursor}
Cursor,
``Composer / Agent mode,''
product documentation, 2025--2026.

\bibitem{icecoder-spec}
iceCoder Project,
``Dual-Mode Adaptive Runtime Supervisor (V1.3.7),''
implementation specification,
\url{docs/requirement/双模方案2-finish.md}, 2026.

\bibitem{icecoder-bench}
iceCoder Project,
``Tri-Platform Same-Model Evaluation Rubric (v0.1),''
\url{benchMark/md/三平台同模对比评测与裁判评分体系.md}, 2026.

\end{thebibliography}

\appendix
\section{Artifact availability}
\label{app:artifacts}

\begin{itemize}[nosep]
  \item Source: iceCoder repository (TypeScript, Node.js 18+).
  \item Entry points: \texttt{npx tsx src/cli/index.ts run}, Web UI, WebSocket chat, optional MCP.
  \item Benchmark reports: \texttt{benchMark/reports/multi-file-order-pipeline.md}, \texttt{saga-warehouse-reconciliation-basic.md}.
  \item Reproduction: \texttt{npm test} (1{,}165 cases); configure \texttt{data/config.json} for LLM provider.
\end{itemize}

\end{document}
